% ================================================================
% BCT LETTER 60
% Dark PeVatrons as OHC Hopf Defects
% ================================================================
\documentclass[twocolumn,prl,showpacs,preprintnumbers,amsmath,amssymb]{revtex4-1}
\usepackage{amsmath,amssymb,amsfonts,booktabs,hyperref}
\begin{document}

\title{BCT Superfluid Lattice Model --- Letter 60\\
Dark PeVatrons as OHC Hopf Defects:\\
A BCT Mechanism for Unidentified Ultra-High-Energy\\
Gamma-Ray Sources}

\author{Michel Robert Cabri\'e}
\email{ZeroFreeParameters@gmail.com}
\affiliation{Independent Researcher, Victoria, Australia\\
ORCID: 0009-0007-9561-9859}
\date{March 2026}

\begin{abstract}
The LHAASO collaboration has identified ultra-high-energy gamma-ray
sources above 100\,TeV with no counterpart at any other wavelength ---
so-called ``dark PeVatrons.'' The absence of radio, optical, and X-ray
emission is anomalous for all standard acceleration mechanisms (pulsar
wind nebulae, supernova remnants, stellar bow shocks), which require
magnetic fields producing synchrotron radiation as a necessary
by-product. We propose a BCT explanation: the dark PeVatron is a
macroscopic Hopf topological defect in the Octet-Hopfion Condensate
(OHC) vacuum. Particle acceleration proceeds via the Madelung quantum
pressure gradient --- a charge-independent force that produces no
synchrotron emission by construction. The BCT framework predicts a
spectral index
\begin{equation*}
\Gamma = 2 + \alpha_0\,(r_\mathrm{oct}/r_\mathrm{tet})\,(4/\pi) = 2.017,
\end{equation*}
an acceleration quantum $E_{D_4} = \Lambda_\mathrm{QCD}\times S_{D_4}
= 11.22\,\mathrm{GeV}$ per unit Hopf winding, and zero electromagnetic
counterpart. Zero free parameters.
\end{abstract}

\maketitle

\section{The Dark PeVatron Problem}

The LHAASO (Large High Altitude Air Shower Observatory) collaboration
has detected ultra-high-energy gamma-ray emission above 100\,TeV from
a growing catalogue of galactic sources~\cite{LHAASO21}. Among these,
a subset are classified as ``dark PeVatrons'': sources with significance
above $8\sigma$ in gamma-rays but no identified counterpart at radio,
optical, X-ray, or infrared wavelengths.

The archetypal example is LASSO J2108+5157~\cite{LHAASO22} --- detected
at $9.5\sigma$ in the 25--100\,TeV band and $8.5\sigma$ above 100\,TeV,
point-like at the instrument PSF, and completely dark at all other
wavelengths.

Standard PeVatron mechanisms all require electromagnetic by-products.
The dark PeVatron has none. BCT provides a mechanism in which this is
not merely possible but required.

\section{OHC Hopf Defects}

The BCT vacuum ground state is the Octet-Hopfion Condensate (OHC):
a field configuration $\Psi_2: S^3 \to S^2$ with Hopf invariant
$H\in\mathbb{Z}$, classified by $\pi_3(S^2)=\mathbb{Z}$
\cite{BCT_L45,BCT_JH}.

The OHC is governed by the Gross-Pitaevskii equation:
\begin{equation}
i\hbar\,\partial_t\Psi = \left[-\frac{\hbar^2}{2m}\nabla^2 + g|\Psi|^2\right]\Psi.
\end{equation}

A Hopf defect of winding $H$ is a topological soliton where $\Psi\to 0$
at the core. It is stable against decay: the tunnelling amplitude to
$H=0$ is $\exp(-S_\mathrm{Hopf})\approx 10^{-2312}$~\cite{BCT_L45}.
The radial healing profile satisfies:
\begin{equation}
f'' + \frac{1}{r}f' - \frac{H^2}{r^2}f + (1-f^2)f = 0,
\end{equation}
with $f(r)\to(r/\xi)^{|H|}$ near the core and $f\to 1$ at large $r$.
The inner and outer radii of the acceleration zone are set by the
BCT void geometry:
\begin{equation}
\frac{R_\mathrm{inner}}{R_\mathrm{outer}} = \frac{r_\mathrm{tet}}{r_\mathrm{oct}}
= \frac{1}{1.8430}.
\end{equation}

\section{The Madelung Acceleration Mechanism}

Writing $\Psi = \sqrt{\rho}\,e^{i\theta}$, the Madelung transformation
separates the GP equation into a continuity equation and quantum Euler
equation:
\begin{equation}
m\,\frac{Dv}{Dt} = -\nabla(g\rho) +
\frac{\hbar^2}{2m}\nabla\!\left[\frac{\nabla^2\!\sqrt{\rho}}{\sqrt{\rho}}\right].
\end{equation}

The second term is the quantum pressure (Bohm potential). Near a Hopf
defect core, where $\rho$ varies from 0 to $\rho_0$, this force
dominates and is:
\begin{itemize}
\item \emph{Charge-independent}: acts on protons, neutrons, and
electrons equally. No Lorentz force. No cyclotron motion.
\item \emph{Synchrotron-free}: no magnetic field, no deflection,
no radio or X-ray emission.
\item \emph{Long-range}: extends continuously from $R_\mathrm{inner}$
to $R_\mathrm{outer}$.
\end{itemize}

The only electromagnetic output is gamma-rays from hadronic $pp$
interactions at the defect boundary. This is exactly the observed
phenomenology.

\section{Spectral Index from the Bogoliubov Spectrum}

The Bogoliubov excitation spectrum on the Hopf defect background,
with the $\tanh^{2|H|}$ core profile, gives an injection spectrum
$dN/dE \propto E^{-\Gamma}$ with:
\begin{equation}
\Gamma_\mathrm{BCT} = 2 + \alpha_0\,\frac{r_\mathrm{oct}}{r_\mathrm{tet}}
\cdot\frac{4}{\pi},
\end{equation}
where $\alpha_0 = r_\mathrm{oct}\times r_\mathrm{tet}/\pi$ is the BCT
crystal coupling constant. Computing:
\begin{align}
\alpha_0 &= 0.0074081, \\
r_\mathrm{oct}/r_\mathrm{tet} &= 1.8430, \\
\Gamma_\mathrm{BCT} &= 2 + 0.0074081\times 1.8430\times(4/\pi) \notag\\
&= \mathbf{2.017}.
\end{align}

Observed PeVatron spectral indices cluster in the range $2.0$--$2.4$.
BCT places the dark PeVatron at the soft end, consistent with
point-like hard-spectrum sources. Zero free parameters.

\section{The D4 Instanton Acceleration Quantum}

The D4 instanton action $S_{D_4} = \pi^5/6 = 51.003$ sets a natural
energy quantum for Hopf defect excitations:
\begin{equation}
E_{D_4} = \Lambda_\mathrm{QCD}\times S_{D_4} = 220\,\mathrm{MeV}
\times 51.003 = 11.22\,\mathrm{GeV}.
\end{equation}

A defect of winding $H$ carries total acceleration potential:
\begin{equation}
E_\mathrm{total}(H) = H\times E_{D_4}\times
\left(1 - \frac{r_\mathrm{tet}}{r_\mathrm{oct}}\right)
= H\times 11.22\,\mathrm{GeV}\times 0.4574.
\end{equation}

For $H\sim 2\times 10^4$, $E_\mathrm{total}\sim 100$\,TeV. These
winding numbers are physically reasonable for galactic-scale OHC
structures: a defect threading 1\,pc at Planck lattice spacing
contains $\sim 10^{100}$ sites, and $H\sim 10^4$ represents fractional
winding density $\sim 10^{-96}$.

\section{Open Calculation: $H_\mathrm{opt}$ and $E_\mathrm{max}$}

The quantitative prediction of $E_\mathrm{max}$ requires $H_\mathrm{opt}$
from the OHC Green's function $G_H(r,r')$ on the Hopf defect background.
The full calculation minimises:
\begin{equation}
E(H) = H^2\,E_\mathrm{core} - H\,E_\mathrm{boundary},
\end{equation}
subject to galactic boundary conditions, where $E_\mathrm{core}$ and
$E_\mathrm{boundary}$ are determined by the Bogoliubov operator
$\mathcal{L}_H$ on the curved condensate background.

This is a priority open calculation in the BCT programme, directly
analogous to the OHC density-of-states gap in the black hole entropy
calculation (Letter 51~\cite{BCT_L51}). The mechanism and spectral
index are established; the quantitative $E_\mathrm{max}$ awaits the
full Bogoliubov treatment.

\section{Predictions Summary}

\begin{table}[h]
\centering
\caption{BCT predictions for dark PeVatrons.}
\begin{tabular}{lll}
\toprule
Observable & BCT Prediction & Status \\
\midrule
Spectral index $\Gamma$ & $2.017$ (0 params) & Falsifiable \\
EM counterpart & None & Confirmed \checkmark \\
Accel.\ quantum & $11.22$\,GeV & Spectral break \\
Max energy $E_\mathrm{max}$ & $H_\mathrm{opt}\times 5.13$\,GeV & Open \\
Radio polarisation & Zero & Testable \\
\bottomrule
\end{tabular}
\end{table}

\section{Conclusion}

Dark PeVatrons are naturally explained in BCT as macroscopic Hopf
topological defects in the OHC vacuum. The Madelung quantum pressure
gradient accelerates particles without magnetic deflection, producing
gamma-rays while remaining dark at all wavelengths. The spectral index
\begin{equation}
\Gamma = 2 + \alpha_0\,(r_\mathrm{oct}/r_\mathrm{tet})\,(4/\pi)
= \mathbf{2.017}
\end{equation}
follows from BCT geometry with zero free parameters.

\begin{thebibliography}{99}
\bibitem{LHAASO21} LHAASO Collaboration, Science \textbf{373}, 425 (2021).
\bibitem{LHAASO22} LHAASO Collaboration, Phys.\ Rev.\ Lett.\ \textbf{128}, 051102 (2022).
\bibitem{BCT_L45} M.~R.~Cabri\'e, BCT Letter 45 (2026), Zenodo.
\bibitem{BCT_JH}  M.~R.~Cabri\'e, BCT Appendix JH (2026), Zenodo:18975018.
\bibitem{BCT_L51} M.~R.~Cabri\'e, BCT Letter 51 (2026), Zenodo.
\bibitem{BCT_PRL} M.~R.~Cabri\'e, BCT PRL Letter (2026), Zenodo:18884415.
\end{thebibliography}

\end{document}
